\documentclass[11pt]{article}
\usepackage[a4paper,margin=26mm]{geometry}
\usepackage{amsmath,amssymb,mathtools}
\usepackage{unicode-math}
\setmainfont{Latin Modern Roman}
\setmathfont{Latin Modern Math}
\usepackage{xurl,hyperref}
\hypersetup{hidelinks,pdftitle={Problem statement}}
\newcommand{\sref}[2]{\hyperlink{source:#1:#2}{[S#2]}}
\newcommand{\cataloguescope}[1]{#1}
\begin{document}
\section{Hyperbolic geometry — Question 10.1}\label{statement}

\subsection{Definitions and mathematical statement}

\paragraph{Shared notation.}$\mathbb N=\{1,2,\ldots\}$, and $\mathbb Z,\mathbb Q,\mathbb R,\mathbb C$ denote the integers, rationals, reals and complex numbers. Any other conventions are specified locally.


A codimension-one foliation $\mathcal F$ decomposes a $3$-manifold $M$ into surface leaves, locally as the horizontal slices of a product chart. In the terminology of the source, it is taut if a closed curve transverse to $\mathcal F$ meets every leaf. A minimal foliation means that every leaf is dense; a minimal surface means a surface with vanishing mean curvature for the specified Riemannian metric. These two uses of minimal are distinct. The leaf space $L_{\mathcal F}=\widetilde M/\widetilde{\mathcal F}$ has its quotient topology, and can be non-Hausdorff. A taut foliation is $\mathbb R$-covered if this space is homeomorphic to $\mathbb R$. Its deck action is the holonomy representation $\rho_H:\pi_1(M)\longrightarrow\operatorname{Homeo}(L_{\mathcal F})$. Branching means nonseparated leaves in this space; one-sided branching occurs in only one transverse direction, and two-sided branching in both.
\paragraph{Statement recorded in the supplied source.}
Let $M$ be hyperbolic and $\mathcal F$ a taut foliation with two-sided branching. Identify $\widetilde M$ with $\mathbb H^3$. Must there exist a lifted leaf $\lambda$ such that each side of $\mathbb H^3\setminus\lambda$ contains an open hyperbolic halfspace? A halfspace is a component of the complement of a totally geodesic plane; a leaf with this property is asymptotically separated. \sref{statement}{2}
\subsection{Short English statement}
Must a taut foliation with branching in both directions have a lifted leaf leaving an open hyperbolic halfspace on each side?
\subsection{Sources}
\begin{description}
\item[\hypertarget{source:statement:1}{S1}] ULAM.AI and the UnsolvedMath Contributors, \emph{UnsolvedMath: A Curated Collection of Open Mathematics Problems}, catalogue entry. CC BY 4.0. \url{https://huggingface.co/datasets/ulamai/UnsolvedMath}.
\item[\hypertarget{source:statement:2}{S2}] Danny Calegari, \emph{Problems in foliations and laminations of 3-manifolds}, version 0.78 (2002), Question 10.1; Definitions 1.1--1.10 and local definitions and remarks. \url{https://arxiv.org/abs/math/0209081v1}.
\end{description}
\label{end:statement}
\end{document}
